% Auto-converted from paper2_full_draft.md §2.
% Wording preserved; no claims added.
\section{Introduction}
Vulnerability triage at scale relies on signals that are easy to share: the \textbf{Exploit Prediction Scoring System (EPSS)}~\cite{Jacobs2021EPSS,Jacobs2023EnhancingEPSS}, the \textbf{CISA Known Exploited Vulnerabilities (KEV) catalog}, \textbf{CVSS v3.1 / v4} severity, and local context such as asset criticality, network exposure, and remediation cost. Practitioners frequently combine these signals with hand-tuned or fitted weights. The fitted-weight path is attractive because calibration promises an estimator of risk grounded in observed exploitation outcomes. It is also potentially fragile: per-feature weight estimation requires \emph{enough labelled positives} — not enough labelled \emph{pairs} (the calibration unit is the CVE, not the host-CVE pair) — and the available public exploit signals are sparse on time-bounded windows.

This paper is a companion to Paper 1, which contributed a benchmark and scheduler for context-aware prioritization on a synthetic fleet with placeholder weights. Paper 1 made no calibration or model-quality claim. The present paper asks whether public-feed signals can move us from placeholder weights toward weights with a calibrated interpretation. The honest answer is no, at the slice we examined. The negative finding is not a bug; it is a data-density boundary that needs to be acknowledged before public-feed calibration becomes a default move. The unit of analysis for per-feature weight learning is the distinct CVE-level positive, not the host-CVE record.

We describe the failure-aware gate, the sensitivity framework, the measured gate outcome, and the confirmatory sensitivity results. We position the work against EPSS / KEV / CVSS / SSVC and against recent capacity-constrained prioritization, EPSS-stability, and context-feature-ablation studies. We do not claim that any fixed prior matches or exceeds EPSS in any general sense. We do not claim deployment readiness. We do claim that the gate methodology is reproducible and that the negative feasibility result is informative.
